\scanpage{228}
\margin{- 192 - \qquad IX}
14.2.\quad\uline{Relations de l'homomorphisme} $i:\underline M_K\longrightarrow G_K$
\uline{avec la monodromie}. Soit $n$ un entier $>0$. On déduit alors de la suite exacte (14.1.7), via la
suite exacte des $\underline{\operatorname{Ext}}^i(\mathbb Z/n\mathbb Z,-)$ sur le site des espaces rigide-analytiques
au-dessus de $K$, une suite exacte canonique (compte tenu du fait que
$\underline M_K^{\mathrm{an}}$ est sans torsion et $G_K^{\mathrm{an}}$ est $n$-divisible) :
\[
0\longrightarrow{}_nG_K^{\mathrm{an}}\longrightarrow{}_nA_K^{\mathrm{an}}
\longrightarrow\underline M_{K,n}^{\mathrm{an}}\longrightarrow0,
\]
où $\underline M_{K,n}\overset{\mathrm{dfn}}=\underline M_K/n\underline M_K$.
Comme tous les termes qui interviennent dans cette suite exacte proviennent de schémas en groupes finis sur $K$,
cette suite exacte peut se réécrire comme provenant d'une suite exacte
\[
(14.2.1)\qquad0\longrightarrow{}_nG_K\longrightarrow{}_nA_K
\longrightarrow\underline M_{K,n}\longrightarrow0.
\]
Pour $n$ variable, ces suites exactes se comportent comme il convient,
et on trouve, par passage à la limite sur les puissances d'un nombre premier $\ell$, une suite exacte analogue
\[
(14.2.2)\qquad0\longrightarrow T_\ell(G_K)\longrightarrow T_\ell(A_K)
\longrightarrow\underline M_K\otimes_{\mathbb Z}\mathbb Z_\ell\longrightarrow0.
\]
Or on a déjà une telle suite exacte par voie purement algébrique, grâce
à la première formule (8.1.2) et à (8.1.7), qui donnent
\[
(14.2.3)\qquad T_\ell(G_K)\xrightarrow{\sim}T_\ell(A_K)^{\mathrm{ef}}\subset T_\ell(A_K),
\quad T_\ell(A_K)/T_\ell(A_K)^{\mathrm{ef}}\simeq\underline M_K\otimes_{\mathbb Z}\mathbb Z_\ell.
\]
Ceci posé, je dis que \uline{les homomorphismes} (14.2.3)
\uline{provenant de la théorie algébrique ne sont autres que ceux déduits de la suite exacte} (14.2.2)
\uline{provenant de la théorie transcendante}. Pour le premier homomorphisme,
correspondant au premier homomorphisme de (14.2.2), notre assertion
résulte aussitôt, par réduction à la situation semi-stable, de la
\margin{510}
\margin{510}
\margin{227}
